% minilecture07.tex: Tutorial 7 mini-lecture: Exoplanets and Habitable Worlds
% Planetary Systems, Kapteyn Institute, University of Groningen
% A 5-10 minute recap, presented by the TAs at the start of Tutorial 7,
% of the Lecture 13-14 concepts that Worksheet 7 trains.

\documentclass[aspectratio=169, 11pt]{beamer}

% ── Theme and macros (shared with the lecture decks) ────────
\usepackage{../../slides/common/beamerthemeIPS}
% macros.tex: Shared math macros for IPS lecture slides
% Mirrors definitions in book/_config.yml for consistency
% Uses \providecommand to avoid conflicts with unicode-math

% ── Derivatives ─────────────────────────────────────────────
\providecommand{\dv}[2]{\frac{\mathrm{d} #1}{\mathrm{d} #2}}
\providecommand{\dd}{}\renewcommand{\dd}{\mathrm{d}}
\providecommand{\pdv}[2]{\frac{\partial #1}{\partial #2}}

% ── Solar quantities ────────────────────────────────────────
\newcommand{\Msun}{M_{\odot}}
\newcommand{\Rsun}{R_{\odot}}
\newcommand{\Lsun}{L_{\odot}}

% ── Planetary quantities ────────────────────────────────────
\newcommand{\Mearth}{M_{\oplus}}
\newcommand{\Rearth}{R_{\oplus}}
\newcommand{\Mjup}{M_{\mathrm{J}}}
\newcommand{\Rjup}{R_{\mathrm{J}}}

% ── Constants ───────────────────────────────────────────────
\newcommand{\kB}{k_{\mathrm{B}}}

% ── Media links ─────────────────────────────────────────────
% Clickable button that opens the animated or video version of a
% figure, e.g. \mediabutton{https://...gif}{Play animation}
\providecommand{\mediabutton}[2]{\href{#1}{\beamergotobutton{#2}}}

\graphicspath{{figures/}{../../slides/lecture13/figures/}{../../slides/lecture14/figures/}{../../slides/common/}}

% ── Metadata ────────────────────────────────────────────────
\title[T7 mini-lecture: Exoplanets and Habitable Worlds]{Tutorial 7 Mini-Lecture:\\Exoplanets and Habitable Worlds}
\subtitle{The concepts behind Worksheet 7 \textbullet\ Lectures 13--14}
\author{}
\institute{Kapteyn Astronomical Institute\\University of Groningen}
\date{2026}
\titlebgcredit{JWST/NIRSpec transit of WASP-39 b. Credit: Alderson et al.\ (2023)}
\titlebgopacity{0.60}

% ════════════════════════════════════════════════════════════
\begin{document}

% ── Title slide ─────────────────────────────────────────────
\begin{frame}[plain,noframenumbering]
  \titlepage
\end{frame}

% ── Roadmap ─────────────────────────────────────────────────
\begin{frame}{What Worksheet 7 trains}
  Five problems, all built on Lectures 13--14. The physics you need, per problem:
  \vskip4pt\relax
  \begin{enumerate}
    \item \textbf{Transits}: the depth $(R_p/R_\star)^2$, the geometric probability $R_\star/a$, why the catalogue is biased.
    \item \textbf{Radial velocities}: the semi-amplitude $K_\star$, from 51 Peg b to the Earth analogue, mass plus radius gives density.
    \item \textbf{Transmission spectroscopy}: the scale height, the $2n_H H/R_p$ modulation, flat spectra and how eclipses break them.
    \item \textbf{The habitable zone}: equilibrium temperature, the Kopparapu flux limits, scaling to K and M dwarfs.
    \item \textbf{Direct imaging and Drake}: angular separations and contrast, the log-uniform Drake toy model.
  \end{enumerate}
  \vskip2pt\relax
  \keyresult{Every problem has the shape and level of one exam question. All parts are analytical; no computer is needed.}
\end{frame}

% ── P1 ──────────────────────────────────────────────────────
\begin{frame}{Transit depths and geometry}
  \small
  \begin{columns}[T]
    \column{0.52\textwidth}
    Depth and probability from Lecture 13:
    \[
      \boxed{\delta = \left(\frac{R_p}{R_\star}\right)^2} \qquad p \approx \frac{R_\star}{a}
    \]
    \begin{itemize}
      \item A hot Jupiter dims a Sun-like star by $\sim 1.5\%$; an Earth analogue by only $\sim 84$ ppm.
      \item At $0.05$ AU the transit probability is $\sim 9\%$; at $1$ AU it drops to $\sim 0.5\%$.
    \end{itemize}

    \column{0.44\textwidth}
    \begin{block}{The catalogue is not the population}
      Close-in planets transit more often, and small stars give deeper transits, so the raw catalogue over-represents both. Occurrence rates follow only after the selection bias is modelled out.
    \end{block}
  \end{columns}
  \keyresult{Two small numbers, $(R_p/R_\star)^2$ and $R_\star/a$, decide what the transit method can see.}
\end{frame}

% ── P2 ──────────────────────────────────────────────────────
\begin{frame}{Radial-velocity masses and densities}
  \small
  \begin{columns}[T]
    \column{0.52\textwidth}
    The stellar reflex semi-amplitude of Lecture 13:
    \[
      \boxed{K_\star = \left(\frac{2\pi G}{P}\right)^{1/3}\frac{m_p \sin i}{M_\star^{2/3}}}
    \]
    \begin{itemize}
      \item 51 Peg b: $K_\star \approx 63\ \mathrm{m\,s^{-1}}$, six times the ELODIE precision.
      \item An Earth analogue: $\approx 9\ \mathrm{cm\,s^{-1}}$, at the ESPRESSO photon limit but below the activity floor.
    \end{itemize}

    \column{0.44\textwidth}
    \begin{block}{From $m_p \sin i$ to composition}
      RV alone gives a minimum mass. A transit forces $\sin i \approx 1$ and adds $R_p$, so mass and radius combine into a bulk density: K2-18 b at $\sim 2700\ \mathrm{kg\,m^{-3}}$ cannot be bare rock.
    \end{block}
  \end{columns}
  \keyresult{The RV amplitude scales as $m_p P^{-1/3}$; combined with a transit it turns into a density and a first look at composition.}
\end{frame}

% ── P3 ──────────────────────────────────────────────────────
\begin{frame}{Scale heights and transmission spectra}
  \small
  \begin{columns}[T]
    \column{0.52\textwidth}
    The atmospheric lever arm of Lecture 13:
    \[
      \boxed{H = \frac{\kB T}{\mu\,m_u\,g}} \qquad \frac{\Delta\delta}{\delta} \approx \frac{2 n_H H}{R_p}
    \]
    \begin{itemize}
      \item Hot Jupiter: $H \approx 630$ km and $\Delta\delta \approx 1120$ ppm, a signal even pre-JWST instruments could see.
      \item Rocky $\mathrm{CO_2}$ planet: $H \approx 5.8$ km, a factor $\sim 109$ smaller signal.
    \end{itemize}

    \column{0.44\textwidth}
    \begin{block}{When the spectrum is flat}
      An airless rock, a heavy compact atmosphere, and high clouds all look featureless in transmission. The $15\,\mu$m eclipse of TRAPPIST-1 b ($503$ K, matching the $508$ K bare-rock value) is what broke the tie there: a thick atmosphere is ruled out, a thin one not yet.
    \end{block}
  \end{columns}
  \keyresult{Hot and light atmospheres are easy; cold and heavy ones hide, and thermal emission has to break the tie.}
\end{frame}

% ── P4 ──────────────────────────────────────────────────────
\begin{frame}{Equilibrium temperature and the habitable zone}
  \small
  \begin{columns}[T]
    \column{0.52\textwidth}
    Radiation balance and the flux limits of Lecture 14:
    \[
      \boxed{T_\mathrm{eq} = \left[\frac{L_\star(1-A_\mathrm{B})}{16\pi\sigma\epsilon d^2}\right]^{1/4}} \quad d = \sqrt{\frac{L_\star/\Lsun}{S_\mathrm{eff}}}\ \mathrm{AU}
    \]
    \begin{itemize}
      \item Earth: $T_\mathrm{eq} \approx 255$ K; the $33$ K to the observed $288$ K is the greenhouse effect.
      \item The Sun's zone spans $0.97$ to $1.69$ AU; at TRAPPIST-1 it shrinks to a few hundredths of an AU.
    \end{itemize}

    \column{0.44\textwidth}
    \begin{block}{Small stars, mixed blessing}
      M dwarfs offer deep transits and long lifetimes, but their habitable-zone planets are tidally locked, baked during the pre-main-sequence phase, and flared for Gyr. K dwarfs may be the best compromise.
    \end{block}
  \end{columns}
  \keyresult{The habitable zone moves as $\sqrt{L_\star}$ while stellar lifetime grows as $M^{-2.5}$: the trade-offs are quantitative, not rhetorical.}
\end{frame}

% ── P5 ──────────────────────────────────────────────────────
\begin{frame}{Direct imaging and the Drake equation}
  \small
  \begin{columns}[T]
    \column{0.52\textwidth}
    Small angles and honest priors from Lectures 13--14:
    \[
      \boxed{\theta[\mathrm{arcsec}] = \frac{a[\mathrm{AU}]}{d[\mathrm{pc}]}}
    \]
    \begin{itemize}
      \item An Earth analogue at $10$ pc sits at $0.1$ arcsec but at $10^{-10}$ contrast: the contrast is the blocker, not the angle.
      \item The M-dwarf habitable zone collapses to milliarcseconds: imaging and transits divide the sky between stellar types.
    \end{itemize}

    \column{0.44\textwidth}
    \begin{block}{Drake as a framing tool}
      With four factors log-uniform over ten decades, $\log_{10} N$ has mean $-20$ and $\sigma \approx 5.8$: the output restates the priors. ``We are alone in the galaxy'' is fully consistent with current knowledge.
    \end{block}
  \end{columns}
  \keyresult{Direct imaging is a contrast problem; the Drake equation is a map of ignorance, and both statements are quantitative.}
\end{frame}

% ── Closer ──────────────────────────────────────────────────
\begin{frame}{Working the worksheet}
  \begin{itemize}
    \item \textbf{Show your algebra.} The exam asks for the steps, not only the number.
    \item Carry \textbf{4 significant figures} through intermediate steps; round at the end.
    \item Use the \textbf{checkpoints} (``show that \ldots'') to verify your work and to keep moving if a step resists.
    \item Qualitative parts want \textbf{2 to 3 sentences} of physics, not an essay.
    \item Most parts close by asking what the number means: that interpretation is graded, not only the arithmetic.
  \end{itemize}
  \vskip8pt\relax
  \keyresult{Worksheet and full solutions: \href{https://ips.formingworlds.space/worksheets.html}{\texttt{ips.formingworlds.space/worksheets.html}}. We discuss questions, not presentations of the solutions.}
\end{frame}

\end{document}
